% Source: sources/returns/S6_04_MAPS_MANIFOLDS/sections/manifolds.tex:1801–2154
\paragraph{Theorem D.2.10 (Proved: unrestricted orders in both Yang families).}
Fix a positive cusp index $R$ in Yang's Type B or Type C period
family, with its canonical cusp collar. Use the following constants:
\[
 \begin{array}{c|ccc|cc}
 \text{type}&c&d&a=d/c&m_1&m_2\\\hline
 B&2&4&2&4k&6l\\
 C&3&3&1&3k&6l
 \end{array}
 \qquad k,l\geq1.
\]
In the marking $\Lambda=\mathbb Z\langle\gamma,u,w,\delta\rangle$,
put $\ell_1=\gamma-R(u+w)$ and choose the invariant finite
translation vectors
\[
 v_1=n_1\ell_1+r_1\delta,\qquad
 v_2=n_2\gamma+r_2\delta,\qquad
 \gcd(n_j,m_j)=1,
 \qquad n_j,r_j\in\mathbb Z.
\]
These data give compact smooth complex threefolds for every displayed
integer choice. Define
\[
 P=m_2n_1+m_1n_2,\qquad h=\gcd(R,ak,P)>0.
\]
If $P\ne0$, set
$G=\mathbb Z/h\oplus\mathbb Z/(R|P|/h)$, omitting trivial factors.
Then
\[
 \begin{array}{c|ccccccc}
 i&0&1&2&3&4&5&6\\\hline
 H_i(X;\mathbb Z)&\mathbb Z&\mathbb Z/|P|&
 \mathbb Z^{R-1}\oplus G&G&
 \mathbb Z^{R-1}\oplus\mathbb Z/|P|&0&\mathbb Z
 \end{array}
\]
and $\pi_1(X)=\mathbb Z/|P|$.
If $P=0$, necessarily $m_1=m_2$ and $n_2=-n_1$; in this case
$h=\gcd(R,ak)$, $\pi_1(X)=\mathbb Z$, and
\[
 \begin{array}{c|ccccccc}
 i&0&1&2&3&4&5&6\\\hline
 H_i(X;\mathbb Z)&\mathbb Z&\mathbb Z&
 \mathbb Z^{R+1}\oplus\mathbb Z/h&
 \mathbb Z^2\oplus\mathbb Z/h&\mathbb Z^{R+1}&
 \mathbb Z&\mathbb Z.
 \end{array}
\]
In both cases $\chi(X)=2R$; the displayed homology is independent
of $r_1,r_2$. No member is simply connected. More generally, with
these two finite multiplicities, this last obstruction holds for
every free affine choice of the two logarithmic fillings, even if
the sufficient coprimality condition above is dropped.

For $k=l=1$, $(n_1,n_2)=(1,-1)$ and $r_1=r_2=0$, this independently
recovers Yang's original groups: Type B has fundamental group
$\mathbb Z/2$ and middle torsion $\mathbb Z/R\oplus\mathbb Z/2$;
Type C has fundamental group $\mathbb Z/3$ and middle torsion
$\mathbb Z/(3R)$.
\begin{proof}
\emph{Analytic construction and exact marking.}
The actual dual deck matrices of the primary source are
\[
 A_1=\begin{pmatrix}
 1&0&0&0\\0&0&1&0\\-cR&-1&2-c&0\\-R&0&-1&1
 \end{pmatrix},\qquad
 A_2=\begin{pmatrix}
 1&0&0&0\\0&1&-1&0\\0&1&0&0\\0&0&0&1
 \end{pmatrix}.
\]
Their orders are $d,6$, and their integral fixed lattices are
$\langle\ell_1,\delta\rangle$ and
$\langle\gamma,\delta\rangle$, respectively.
Pull the good-reduction root discs back by degrees $k,l$.
On the resulting root family the affine action is
\[
 (s,x)\longmapsto(\zeta_{m_j}^{-1}s,A_jx+v_j/m_j).
\]
The holomorphic period formula and the logarithmic collar maps of
\cite[Lemma 4.1, Lemma 4.3]{manYang26} apply unchanged: in complex
coordinates the translation is the holomorphic period section
$\Pi(\zeta_{m_j}^{-1}s)v_j/m_j$.
The invariant functional $\chi=\gamma^*$ changes by $n_j/m_j$.
Applying it to a purported fixed point of a nontrivial power
forces $m_j\mid qn_j$, impossible for $0<q<m_j$.
The actions are therefore free, their quotients are proper, and
the central divisors have multiplicities $m_j$.
The collar isomorphism adds or subtracts
$(\log s)/(2\pi i)\,\Pi(s)v_j$; changing the logarithm adds an
integral period. This proves existence with the stated marking
and leaves the toroidal cusp unchanged. Compactness follows from
the same three compact tubes and compact base core as in the
source's global gluing argument.

Write $g_j$ for the transferred clockwise meridian. The finite
boundary map imposes exactly $g_j^{m_j}=v_j$, by
\cite[Proposition 4.4]{manYang26}; this can also be checked from
the path $q\mapsto qv_j/m_j$ in the central affine quotient.
The cusp kills $w,\delta$ and $g_1g_2$.
Since $A_1w=u+(2-c)w-\delta$, their normal closure kills $u$.
The remaining generators $\gamma,g_1$ commute and satisfy
\[
 m_1g_1=n_1\gamma,\qquad -m_2g_1=n_2\gamma.
\]
The entries of this relation matrix have gcd one, and its
determinant has absolute value $|P|$. This proves the asserted
fundamental groups, including the rank-one group when $P=0$.
For arbitrary free affine fillings one can instead kill the
entire fibre subgroup: the resulting quotient is
\[
 \langle g_1,g_2\mid g_1g_2=1,
                  g_1^{m_1}=g_2^{m_2}=1\rangle
 =\mathbb Z/\gcd(m_1,m_2).
\]
It is nontrivial, because its order is divisible by two in Type B
and by three in Type C. This proves the stronger obstruction.

\emph{Integral finite-quotient lattices.}
Before the cusp is attached, let $X^\circ$ be the union of the two
finite pieces. Each retracts to its reduced quotient surface $S_j$,
and their intersection retracts to a regular torus $F$.
The first linear coinvariants satisfy $w=-u$, $\delta=cu$;
the second kill $u,w$. Thus
\[
 \begin{split}
 H_1(S_1)&=\langle\gamma,u,g_1\mid
             m_1g_1=n_1\gamma+cr_1u\rangle,\\
 H_1(S_2)&=\langle\gamma,\delta,g_2\mid
             m_2g_2=n_2\gamma+r_2\delta\rangle.
 \end{split}
\]
Each relation is primitive by $\gcd(n_j,m_j)=1$; both groups are
$\mathbb Z^2$. Euler characteristic zero, duality and the universal
coefficient theorem give $H^2(S_j)=\mathbb Z^2$ with a unimodular
intersection pairing.

Use the ordered regular-fibre cohomology basis
\[
 (\gamma^*u^*,\gamma^*w^*,\gamma^*\delta^*,
            u^*w^*,u^*\delta^*,w^*\delta^*).
\]
Solving the integral fixed equations gives the first invariant
lattice with basis
\[
 \omega_1=-R\gamma^*u^*+R\gamma^*w^*+u^*w^*,\qquad
 \omega_2=\gamma^*u^*-\gamma^*w^*+c\gamma^*\delta^*.
\]
Both squares vanish and $\omega_1\omega_2=c\,\mathrm{vol}_F$.
The second invariant lattice has basis
$u^*w^*,\gamma^*\delta^*$ and pairing $U$.
If $L_j=p_j^*H^2(S_j)$, transfer gives injectivity and the
discriminant formula gives the exact indices
\[
 [H^2(F)^{A_1}:L_1]=m_1/c=ak,
 \qquad [H^2(F)^{A_2}:L_2]=m_2.
\]
Indeed, the quotient pairing is unimodular and pullback multiplies
it by the covering degree $m_j$.

The square $[q\delta+t v_j/m_j]$ closes to a two-torus in $S_j$.
For a pulled-back integral class $\alpha$ its evaluation is
$\alpha(\delta,v_j)/m_j\in\mathbb Z$.
For $\alpha=A\omega_1+B\omega_2$ this is
$-cn_1B/m_1$; coprimality forces $B\in ak\mathbb Z$.
At the second fibre it forces the coefficient of
$\gamma^*\delta^*$ to lie in $m_2\mathbb Z$.
These candidate sublattices have exactly the required indices,
so they are the entire pullback lattices, without a hidden
saturation factor. Consequently their integral basis matrices are
\[
 P_1=\begin{pmatrix}
 -R&ak\\R&-ak\\0&m_1\\1&0\\0&0\\0&0
 \end{pmatrix},\qquad
 P_2=\begin{pmatrix}
 0&0\\0&0\\0&-m_2\\1&0\\0&0\\0&0
 \end{pmatrix}.
\]
Choose the homology bases dual to these integral cohomology
bases. In particular, when $|n_j|>1$ the finite $\delta$ sweep is
not declared to be a basis vector. Its actual coordinates are
$(0,-n_1)$ and $(0,n_2)$ in the respective surfaces. This is the
integral distinction needed for arbitrary twists.
The degree-two Mayer--Vietoris map is therefore
\[
 \alpha_2=\begin{pmatrix}
 -R&R&0&1&0&0\\ak&-ak&m_1&0&0&0\\
 0&0&0&-1&0&0\\0&0&m_2&0&0&0
 \end{pmatrix}.
\]
The degree-one kernel is exactly the primitive line
$\ker\alpha_1=\mathbb Z(u+w)$: the second inclusion kills only
the $u,w$ directions of $\Lambda$, and the first then forces their
coefficients to agree. The extra relation involving $g_j$ cannot
create another kernel vector from $\Lambda$, because its
$g_j$ coefficient is the nonzero integer $m_j$.

\emph{The two cusp relations are primitive.}
Let $N$ be the unchanged cusp tube, $M=\partial N$, and
$P_0=A_1A_2$ its clockwise forward transport.
Direct exterior-square calculation gives rank two and a unit
$2\times2$ minor of $\bigwedge^2P_0-I$: use columns $(1,4)$
and rows $(2,6)$ in Type B or $(3,6)$ in Type C.
In both cases the determinant is $-1$, for every $R$.
Thus its cokernel is $\mathbb Z^4$. The Wang sequence then makes
$H_2(M)$ free, and the two swept tori $\xi_w,\xi_\delta$ project
to the integral basis $w,\delta$ of $\ker(P_0-I)$.
They generate a primitive direct summand.

For completeness, the normalization calculation used here is
integral. There are $R$ normalized $dP_6$ components, $3R$ double
curves and $2R$ triple points. With $\epsilon=1$ in Type B and
$\epsilon=2$ in Type C, component labels are
$j=a-\epsilon b\pmod R$. If a component class has coordinates
$(A_j,b_{1j},b_{2j},b_{3j})$, its three restriction differences are
\[
 \begin{split}
 X_j&=b_{1j}-A_{j+1}+b_{2,j+1}+b_{3,j+1},\\
 Y_j&=A_j-b_{1j}-b_{2j}-b_{3,j-\epsilon},\\
 Z_j&=b_{2j}-A_{j-\epsilon-1}
                  +b_{1,j-\epsilon-1}+b_{3,j-\epsilon-1}.
 \end{split}
\]
Its image is exactly the kernel of the surjective map
$(X,Y,Z)\mapsto(\sum(X+Y),\sum(Y+Z))$ to $\mathbb Z^2$.
To verify equality, put $T=S^{-1}$ for cyclic shift and use $Y$
to eliminate $A$. Solve
$Z+T^{\epsilon+1}Y=(I-T)q$, possible precisely when the second
sum vanishes; set $b_2=q$, $b_3=-T^{-\epsilon}q$.
The identity
$(I-T^{\epsilon+1})-T(I-T^\epsilon)=I-T$ solves that equation.
The other one is an integer cyclic difference equation for $b_1$
with zero coordinate sum, and is therefore soluble.
This proves saturation for all $R$, including self-identifications
when $R=1$. The normalization spectral sequence, whose other
rows are the dual two-torus and $\mathbb Z^R$ in degree four,
now gives $H^3(N)=\mathbb Z^2$, $H^4(N)=\mathbb Z^R$ and
$\chi(N)=2R$; this checks \cite[Theorem 7.8]{manYang26} directly.

The angular torus action on $N$ sweeps the extended zero-section
disc into integral solid-torus chains bounding $\xi_w,\xi_\delta$.
Poincar\'e--Lefschetz duality gives
$H_3(N,M)=H^3(N)=\mathbb Z^2$, so the kernel of
$H_2(M)\to H_2(N)$ has rank at most two. It contains the primitive
rank-two summand just constructed; a finite enlargement would
embed in the free quotient of $H_2(M)$ by that summand.
Hence
\[
 K_2=\ker(H_2(M)\to H_2(N))
       =\mathbb Z\xi_w\oplus\mathbb Z\xi_\delta.
\]
This establishes the complete kernel over $\mathbb Z$.

The two finite forward transports are $P_j=A_j^{-1}$.
They send $w$ first to $-u$ and then back to $w$.
Cutting the clockwise cusp sweep into the two cylinders gives
connecting class $-(u+w)$ in the convention
$\alpha=(i_{1*},-i_{2*})$, exactly as the explicit cylinder
boundaries in \cite[Lemma 7.10, Appendix B.4]{manYang26}.
It is a primitive generator of $\ker\alpha_1$, irrespective
of the finite translation vectors. Sweeping $\delta$ over the
pair of pants gives
$j_*\xi_\delta=-s_{1,\delta}-s_{2,\delta}$.
The evaluations already computed show that $-j_*\xi_\delta$
has the column $b=(0,-n_1,0,n_2)^{\mathsf T}$ in the integral
dual bases. Therefore
\[
 H_2(X^\circ)/j_*K_2
    =\operatorname{coker}[\alpha_2\mid b].
\]
Eliminate the unit relation from the $uw$ column. The remaining
three-generator presentation is
\[
 C=\begin{pmatrix}
 R&0&0\\-ak&m_1&-n_1\\0&m_2&n_2
 \end{pmatrix},\qquad \det C=RP.
\]
Its entry gcd is one. Its nonzero $2\times2$ minors, up to signs,
are
\[
 Rm_1,\ Rn_1,\ Rm_2,\ Rn_2,\ akm_2,\ akn_2,\ P.
\]
Their gcd is exactly $h=\gcd(R,ak,P)$, by the two coprimality
conditions. Thus for $P\ne0$ its Smith factors are
$(1,h,R|P|/h)$; for $P=0$ its rank is two and its factors are
$(1,h,0)$. The latter cokernel is $\mathbb Z\oplus\mathbb Z/h$.

\emph{The relative term, including $P=0$.}
Excision and duality give $H_2(N,M)=H^4(N)=\mathbb Z^R$.
Before the cusp is attached, the joint monodromy images kill
$\ker\chi$ in the abelianization. Hence
\[
 H_1(X^\circ)=\left\langle\gamma,g_1,g_2\ \middle|\
 m_1g_1=n_1\gamma,\quad m_2g_2=n_2\gamma\right\rangle
       =\mathbb Z\oplus\mathbb Z/d_0,
 \quad d_0=\gcd(m_1,m_2).
\]
The equality follows integrally: the relation matrix has entry
gcd one and maximal-minor gcd $d_0$.
The full image of the relative boundary map is
$\langle a_0\rangle$, where $a_0=-(g_1+g_2)$.
Indeed, the cusp boundary kernel is generated by its meridian
and its vanishing lattice; the latter is already zero in
$H_1(X^\circ)$, and the extended section disc supplies $a_0$.

When $P\ne0$, the primitive homomorphism from $H_1(X^\circ)$
to its free quotient is
\[
 \Phi(\gamma,g_1,g_2)
 =\frac1{d_0}(m_1m_2,\ m_2n_1,\ m_1n_2).
\]
The same gcd-of-minors calculation proves primitivity.
It sends $a_0$ to $-P/d_0\ne0$, so $\langle a_0\rangle$
is infinite cyclic. The kernel of the relative boundary map
is free of rank $R-1$. This does not assume that $a_0$ is
primitive in $H_1(X^\circ)$.

When $P=0$, equality of the reduced fractions
$n_1/m_1=-n_2/m_2$ forces $m_1=m_2=m$ and $n_2=-n_1$.
Use $a_0$ and $g_1$ as meridian generators: the relations are
$ma_0=0$ and $mg_1=n_1\gamma$. The second is primitive,
so $a_0$ has exact order $m$. The relative boundary image is
therefore the finite group $\mathbb Z/m$, and its kernel is
free of rank $R$, of index $m$ in $\mathbb Z^R$.
This is the extra free contribution in the degenerate case.

The exact sequence for $(X,X^\circ)$ now reads
\[
 0\longrightarrow\operatorname{coker}C
 \longrightarrow H_2(X)
 \longrightarrow\ker\bigl(H_2(N,M)\to H_1(X^\circ)\bigr)
 \longrightarrow0.
\]
Its right group is free, so it splits. This gives the displayed
degree-two groups in both cases. All pieces other than the cusp
have Euler characteristic zero, so $\chi(X)=2R$.
For $P\ne0$, duality gives $b_1=b_5=0$, $b_2=b_4=R-1$,
and Euler characteristic gives $b_3=0$.
For $P=0$, it gives $b_1=b_5=1$, $b_2=b_4=R+1$ and then
$b_3=2$. In either case the universal coefficient theorem and
duality identify the degree-three torsion with the degree-two
torsion and the degree-four torsion with the degree-one torsion.
This proves all the integral homology assertions separately
in the two cases.

The script \texttt{code/manifolds\_yang.py} verifies the symbolic
monodromy and primitive-minor identities, 32 full cusp
normalization Smith forms, 162 finite-quotient and sweep
checks, and 2592 full augmented $4\times7$ Smith forms,
including 72 with $P=0$. Its executable also accepts arbitrary
admissible integer inputs. These finite checks accompany the
unbounded integral proof; they impose no hypothesis on it.
The original cases are \cite[Theorem 7.13]{manYang26}; no
priority claim is made for them or for the extensions above.
\end{proof}
